	\begin{UseCase}{CU12.1.1.1.1}{Gestionar Pasos}{
     		El colaborador (\hyperlink{jefe}{Líder de Análisis} o \hyperlink{analista}{Analista}) podrá gestionar los \hyperlink{entidadPaso}{Pasos} propios de cada \hyperlink{entidadTray}{Trayectoria}.\\
     	
	     	Los pasos que describen una trayectoria deben indicar las acciones que realiza el sistema o el actor. Este caso de uso permite al colaborador visualizar en una tabla, el registro de los Pasos que hay en una Trayectoria del caso de uso sobre el cual se está operando, así como solicitar el registro de un paso nuevo, modificar alguno existente (teniendo la posibilidad de actualizar algún dato de registro), o eliminarlo (en caso de que la acción por algún motivo ya no forme parte de la Trayectoria)\\
	     	
	     	Las acciones disponibles para cada Paso dependerán del estado en el que se encuentre el caso de uso donde son referenciadas. 
     	
	}
	\UCitem{Actor}{\hyperlink{jefe}{Líder de análisis}, \hyperlink{analista}{Analista}}
	\UCitem{Propósito}{Proporcionar al actor un mecanismo para llevar el control de pasos que componen a una trayectoria de un caso de uso.}
	\UCitem{Entradas}{Ninguna}
	\UCitem{Salidas}{\begin{itemize}
			\item \cdtRef{proyectoEntidad:claveProyecto}{Clave del proyecto}: Lo obtiene el sistema.
			\item \cdtRef{proyectoEntidad:nombreProyecto}{Nombre del proyecto}: Lo obtiene el sistema.
			\item \cdtRef{moduloEntidad:nombreModulo}{Nombre del Módulo}: Lo obtiene el sistema.
			\item \cdtRef{casoUso:numeroCU}{Número} del caso de uso: Lo obtiene el sistema. 
			\item \cdtRef{casoUso:nombreCU}{Nombre} del caso de uso: Lo obtiene el sistema.
			\item \cdtRef{entidadTray:nombreTray}{Nombre} de la trayectoria: Lo obtiene el sistema.
			\item \cdtRef{entidadPaso}{Paso}: Tabla que muestra el \cdtRef{entidadPaso:numeroPaso}{Número} y \cdtRef{entidadPaso:radaccionPaso}{redacción} de todos los pasos registrados en una trayectoria.
			\item \cdtIdRef{MSG2}{No existe información}: Se muestra en la pantalla \IUref{IU6.1.1.1.1.1}{Gestionar Pasos} cuando no existen pasos registrados.
	\end{itemize}}
	
	\UCitem{Precondiciones}{
		\begin{itemize}
			\item Que el caso de uso se encuentre en estado ''Edición'' o ''Pendiente de corrección''.
			\item Que exista al menos una trayectoria registrada.
		\end{itemize}
	}
	\UCitem{Postcondiciones}{Ninguna}
	\UCitem{Errores}{Ninguno}
	\UCitem{Tipo}{Secundario, extiende del caso de uso \UCref{CU12.1.1}{Gestionar Trayectorias}.}
\end{UseCase}
%--------------------------------------
\begin{UCtrayectoria}
	\UCpaso[\UCactor] Solicita gestionar los pasos de una trayectoria seleccionando el icono \raisebox{-1mm}{\includegraphics[height=11pt]{images/Iconos/pasos}} de la pantalla \IUref{IU6.1.1}{Gestionar Trayectorias}.
	\UCpaso[\UCsist] Obtiene la información de los pasos registrados de la trayectoria seleccionada. \hyperlink{CU12-1-1-1-1:TAA}{[Trayectoria A]}
	\UCpaso[\UCsist] Ordena los pasos en base en el orden en el que van siendo creados.
	\UCpaso[\UCsist] Muestra la información de los pasos en la pantalla \IUref{IU6.1.1.1.1.1}{Gestionar Pasos}. \label{CU12-1-1-1-1-P4}
	\UCpaso[\UCactor] Gestiona las pasos a través de los botones: \IUbutton{Registrar}, \editar y \eliminar. 
\end{UCtrayectoria}		
%--------------------------------------
\hypertarget{CU12-1-1-1-1:TAA}{\textbf{Trayectoria alternativa A}}\\
\noindent \textbf{Condición:} No existen registros de pasos.
\begin{enumerate}
	\UCpaso[\UCsist] Muestra el mensaje \cdtIdRef{MSG2}{No existe información} en la pantalla \IUref{IU6.1.1.1.1.1}{Gestionar Pasos} para indicar que no hay registros de pasos para mostrar.  \label{CU12-1-1-1-1-TA1}
	\UCpaso[\UCactor] Gestiona las pasos a través del botón: \IUbutton{Registrar}. 
	\item[- -] - - {\em {Fin de la trayectoria}}.%
\end{enumerate}
%--------------------------------------
\subsubsection{Puntos de extensión}

\UCExtenssionPoint{El actor requiere registrar un paso de la trayectoria.}{Presionando el botón \IUbutton{Registrar} del paso \ref{CU12-1-1-1-1-P4} de la trayectoria principal o del paso \ref{CU12-1-1-1-1-TA1} de la Trayectoria alternativa A.}{\UCref{CU12.1.1.1.1.1}{Registrar Paso}}
\UCExtenssionPoint{El actor requiere modificar un paso de la trayectoria.}{Presionando el icono \editar del paso \ref{CU12-1-1-1-1-P4} de la trayectoria principal.}{\UCref{CU12.1.1.1.1.2}{Modificar Paso}}
\UCExtenssionPoint{El actor requiere eliminar un paso de la trayectoria.}{Presionando el icono \eliminar del paso \ref{CU12-1-1-1-1-P4} de la trayectoria principal.}{\UCref{CU12.1.1.1.1.3}{Eliminar Paso}}
